% SPDX-License-Identifier: GPL-3.0-or-later
% Copyright (C) 2026 Alessandro Greco
%% Descripción alternativa accesible, solo para EPUB.
\altfig{En el ataque histórico de 2014, un HSDir malicioso codifica el nombre del servicio onion en los tipos de celda RELAY y RELAY\_EARLY enviados al cliente. Un guardia malicioso cómplice lee esos tipos visibles para los repetidores y vincula al cliente con la solicitud de un descriptor, no con una visita demostrada. Las cuatro celdas dibujadas son esquemáticas; TLS protege las cabeceras en los enlaces.}
% Ataque de 2014: señal en el tipo de celda leída por el guardia
\begin{tikzpicture}[font=\sffamily\scriptsize,
  ostile/.style={relay,fill=avviso!20,draw=avviso,align=center,font=\sffamily\scriptsize}]
  \node[client] (c) at (0,0) {cliente};
  \node[ostile] (g) at (3,0) {guardia\\malicioso};
  \node[relay] (m) at (6,0) {intermedio};
  \node[ostile] (h) at (9,0) {HSDir\\malicioso};
  \draw[<-,thick] (c) -- (g);
  \draw[<-,thick] (g) -- (m);
  \draw[<-,thick] (m) -- (h);
  \foreach \x/\t/\c in {7.1/E/avviso!40,7.5/R/white,7.9/R/white,8.3/E/avviso!40} {
    \node[draw,minimum size=3.5mm,inner sep=0pt,fill=\c] at (\x,0.55) {\t};
  }
  \node[align=center,text=torgrigio] at (7.6,1.12) {esquema: cuatro celdas de ejemplo};
  \node[align=center,text=avviso] at (9,-1) {codifica el nombre\\en tipos RELAY (R) y\\RELAY\_EARLY (E)};
  \node[align=center,text=avviso] at (3,-1) {lee los tipos:\\vincula al cliente\\con la solicitud};
  \node[align=center,text=torgrigio] at (6,-1) {reenvía los tipos\\como en el caso de 2014};
\end{tikzpicture}
